\documentclass[10pt,a4paper]{amsart}
\usepackage{amsmath,amssymb,amsthm,graphicx}
% Retained preamble for the Stacks Project chapter snapshot in this directory.
% The chapter's own preamble.tex is a copy of the shared Stacks Project preamble; the block
% below is sliced MECHANICALLY, line for line, from it. Nothing was added, removed or
% redefined; the wrapper class, the xr-hyper/hyperref wiring and the externaldocument list
% are omitted because the assembler supplies the document class.
%
% lines taken from preamble.tex: 3, 5, 9, 11, 12, 156, 157, 158, ... (51 lines total)
%
% package coverage restored from the chapter's own preamble (2026-10-13): xy, xy-2cell

\usepackage[all]{xy}
\xyoption{2cell}
\UseAllTwocells

\newenvironment{reference}{\comment}{\endcomment}
\newenvironment{slogan}{\comment}{\endcomment}
\newenvironment{history}{\comment}{\endcomment}
\theoremstyle{plain}
\newtheorem{theorem}[subsection]{Theorem}
\newtheorem{proposition}[subsection]{Proposition}
\newtheorem{lemma}[subsection]{Lemma}
\theoremstyle{definition}
\newtheorem{definition}[subsection]{Definition}
\newtheorem{example}[subsection]{Example}
\newtheorem{exercise}[subsection]{Exercise}
\newtheorem{situation}[subsection]{Situation}
\theoremstyle{remark}
\newtheorem{remark}[subsection]{Remark}
\newtheorem{remarks}[subsection]{Remarks}
\numberwithin{equation}{subsection}
\def\lim{\mathop{\mathrm{lim}}\nolimits}
\def\colim{\mathop{\mathrm{colim}}\nolimits}
\def\Spec{\mathop{\mathrm{Spec}}}
\def\Hom{\mathop{\mathrm{Hom}}\nolimits}
\def\Ext{\mathop{\mathrm{Ext}}\nolimits}
\def\SheafHom{\mathop{\mathcal{H}\!\mathit{om}}\nolimits}
\def\SheafExt{\mathop{\mathcal{E}\!\mathit{xt}}\nolimits}
\def\Sch{\mathit{Sch}}
\def\Mor{\mathop{\mathrm{Mor}}\nolimits}
\def\Ob{\mathop{\mathrm{Ob}}\nolimits}
\def\Sh{\mathop{\mathit{Sh}}\nolimits}
\def\NL{\mathop{N\!L}\nolimits}
\def\CH{\mathop{\mathrm{CH}}\nolimits}
\def\proetale{{pro\text{-}\acute{e}tale}}
\def\etale{{\acute{e}tale}}
\def\QCoh{\mathit{QCoh}}
\def\Ker{\mathop{\mathrm{Ker}}}
\def\Im{\mathop{\mathrm{Im}}}
\def\Coker{\mathop{\mathrm{Coker}}}
\def\Coim{\mathop{\mathrm{Coim}}}
\DeclareMathSymbol{\boxtimes}{\mathbin}{AMSa}{"02}
\def\QCohstack{\mathcal{QC}\!\mathit{oh}}
\def\Cohstack{\mathcal{C}\!\mathit{oh}}
\def\Spacesstack{\mathcal{S}\!\mathit{paces}}
\def\Quotfunctor{\mathrm{Quot}}
\def\Hilbfunctor{\mathrm{Hilb}}
\def\Curvesstack{\mathcal{C}\!\mathit{urves}}
\def\Polarizedstack{\mathcal{P}\!\mathit{olarized}}
\def\Complexesstack{\mathcal{C}\!\mathit{omplexes}}
\def\Pic{\mathop{\mathrm{Pic}}\nolimits}
\def\Picardstack{\mathcal{P}\!\mathit{ic}}
\def\Picardfunctor{\mathrm{Pic}}
\def\Deformationcategory{\mathcal{D}\!\mathit{ef}}
\usepackage{cleveref}
\crefname{originaldefinition}{original definition}{original definitions}
\Crefname{originaldefinition}{Original definition}{Original definitions}
\makeatletter
\newlabel{definition-strict-normal-crossings}{{0BI9}{source}{}{sourcepointer-1}{}}
\newlabel{definition-strict-normal-crossings@cref}{{[originaldefinition][0][]0BI9}{[1][0][]source}{}{}{}}
\makeatother
\crefname{originalreference}{original reference}{original references}
\Crefname{originalreference}{Original reference}{Original references}
\makeatletter
\newlabel{algebra-lemma-regular-ring-CM}{{00NQ}{source}{}{sourcepointer-2}{}}
\newlabel{algebra-lemma-regular-ring-CM@cref}{{[originalreference][0][]00NQ}{[1][0][]source}{}{}{}}
\makeatother
\crefname{originalreference}{original reference}{original references}
\Crefname{originalreference}{Original reference}{Original references}
\makeatletter
\newlabel{algebra-lemma-regular-mod-x}{{00NU}{source}{}{sourcepointer-3}{}}
\newlabel{algebra-lemma-regular-mod-x@cref}{{[originalreference][0][]00NU}{[1][0][]source}{}{}{}}
\makeatother
\crefname{originalreference}{original reference}{original references}
\Crefname{originalreference}{Original reference}{Original references}
\makeatletter
\newlabel{properties-lemma-normal-Noetherian}{{033M}{source}{}{sourcepointer-4}{}}
\newlabel{properties-lemma-normal-Noetherian@cref}{{[originalreference][0][]033M}{[1][0][]source}{}{}{}}
\makeatother
\crefname{originalreference}{original reference}{original references}
\Crefname{originalreference}{Original reference}{Original references}
\makeatletter
\newlabel{algebra-lemma-CM-dim-formula}{{00NA}{source}{}{sourcepointer-5}{}}
\newlabel{algebra-lemma-CM-dim-formula@cref}{{[originalreference][0][]00NA}{[1][0][]source}{}{}{}}
\makeatother
\title{Strict normal crossings divisors characterized by reducedness and regular intersections of components}
\author{The Stacks Project Authors}
\date{stacks-project:etale.tex:local-snapshot-20261012; GNU FDL 1.2 or later}
\begin{document}
\maketitle
\noindent\textit{Editorial note.} All delivered mathematical text is the author's own, copied verbatim from the pinned Stacks Project chapter file etale.tex; no statement or proof step was rewritten, reordered or removed.\par\medskip
\noindent\textit{Editorial note.} Span recomputation by label pairing (LANE\_BRIEF 10.2): in this chapter the author's proof environment follows the statement environment as a standalone block; the spec covers the complete environment pair (statement span, proof span) as recorded in the pool candidate.\par\medskip
\ifdefined\hypertarget\hypertarget{sourcepointer-1}{}\fi
\noindent\textit{Source reference.} Original definition 0BI9: sources/stacks-chow/tags.txt, source lines 14482--14482. Cross-chapter standard prerequisite cited by the author inside the counted span; it is used as a pointer and its own proof is not part of this item.\par\medskip
\ifdefined\hypertarget\hypertarget{sourcepointer-2}{}\fi
\noindent\textit{Source reference.} Original reference 00NQ: sources/stacks-chow/tags.txt, source lines 868--868. Cross-chapter standard prerequisite cited by the author inside the counted span; it is used as a pointer and its own proof is not part of this item.\par\medskip
\ifdefined\hypertarget\hypertarget{sourcepointer-3}{}\fi
\noindent\textit{Source reference.} Original reference 00NU: sources/stacks-chow/tags.txt, source lines 872--872. Cross-chapter standard prerequisite cited by the author inside the counted span; it is used as a pointer and its own proof is not part of this item.\par\medskip
\ifdefined\hypertarget\hypertarget{sourcepointer-4}{}\fi
\noindent\textit{Source reference.} Original reference 033M: sources/stacks-chow/tags.txt, source lines 4061--4061. Cross-chapter standard prerequisite cited by the author inside the counted span; it is used as a pointer and its own proof is not part of this item.\par\medskip
\ifdefined\hypertarget\hypertarget{sourcepointer-5}{}\fi
\noindent\textit{Source reference.} Original reference 00NA: sources/stacks-chow/tags.txt, source lines 852--852. Cross-chapter standard prerequisite cited by the author inside the counted span; it is used as a pointer and its own proof is not part of this item.\par\medskip
\begin{lemma}
\label{lemma-strict-normal-crossings}
Let $X$ be a locally Noetherian scheme. Let $D \subset X$ be an
effective Cartier divisor. Let $D_i \subset D$, $i \in I$ be its
irreducible components viewed as reduced closed subschemes of $X$.
The following are equivalent
\begin{enumerate}
\item $D$ is a strict normal crossings divisor, and
\item $D$ is reduced, each $D_i$ is an effective Cartier divisor, and
for $J \subset I$ finite the scheme theoretic
intersection $D_J = \bigcap_{j \in J} D_j$ is a
regular scheme each of whose irreducible components has
codimension $|J|$ in $X$.
\end{enumerate}
\end{lemma}

\begin{proof}
Assume $D$ is a strict normal crossings divisor. Pick $p \in D$
and choose a regular system of parameters $x_1, \ldots, x_d \in \mathfrak m_p$
and $1 \leq r \leq d$ as in
Definition \ref{definition-strict-normal-crossings}.
Since $\mathcal{O}_{X, p}/(x_i)$ is a regular local ring
(and in particular a domain) we see that the irreducible components
$D_1, \ldots, D_r$ of $D$ passing through $p$ correspond $1$-to-$1$
to the height one primes $(x_1), \ldots, (x_r)$ of $\mathcal{O}_{X, p}$.
By Algebra, Lemma \ref{algebra-lemma-regular-ring-CM}
we find that the intersections $D_{i_1} \cap \ldots \cap D_{i_s}$
have codimension $s$ in an open neighbourhood of $p$
and that this intersection has a regular local ring at $p$.
Since this holds for all $p \in D$ we conclude that (2) holds.

\medskip\noindent
Assume (2). Let $p \in D$. Since $\mathcal{O}_{X, p}$ is finite
dimensional we see that $p$ can be contained in at most
$\dim(\mathcal{O}_{X, p})$ of the components $D_i$.
Say $p \in D_1, \ldots, D_r$ for some $r \geq 1$.
Let $x_1, \ldots, x_r \in \mathfrak m_p$ be local equations
for $D_1, \ldots, D_r$. Then $x_1$ is a nonzerodivisor in $\mathcal{O}_{X, p}$
and $\mathcal{O}_{X, p}/(x_1) = \mathcal{O}_{D_1, p}$ is regular.
Hence $\mathcal{O}_{X, p}$ is regular, see
Algebra, Lemma \ref{algebra-lemma-regular-mod-x}.
Since $D_1 \cap \ldots \cap D_r$ is a regular (hence normal) scheme
it is a disjoint union of its irreducible components
(Properties, Lemma \ref{properties-lemma-normal-Noetherian}).
Let $Z \subset D_1 \cap \ldots \cap D_r$
be the irreducible component containing $p$.
Then $\mathcal{O}_{Z, p} = \mathcal{O}_{X, p}/(x_1, \ldots, x_r)$
is regular of codimension $r$ (note that since we already know
that $\mathcal{O}_{X, p}$ is regular and hence Cohen-Macaulay,
there is no ambiguity about codimension as the ring is catenary, see
Algebra, Lemmas \ref{algebra-lemma-regular-ring-CM} and
\ref{algebra-lemma-CM-dim-formula}).
Hence $\dim(\mathcal{O}_{Z, p}) = \dim(\mathcal{O}_{X, p}) - r$.
Choose additional $x_{r + 1}, \ldots, x_n \in \mathfrak m_p$
which map to a minimal system of generators of $\mathfrak m_{Z, p}$.
Then $\mathfrak m_p = (x_1, \ldots, x_n)$ by Nakayama's lemma
and we see that $D$ is a normal crossings divisor.
\end{proof}

\end{document}
